As large language models become capable of automating the research lifecycle, evaluating their capacity for genuine empirical discovery is critical. Existing benchmarks often use standard datasets or textbook physics, making it impossible to distinguish between true scientific inference and the regurgitation of memorized training data. We introduce the \methodname, a methodological framework where an AI Scientist is tasked with discovering a counterfactual physical law from synthetic data, allowing us to evaluate its reliance on semantic priors versus empirical evidence. We show that when variables have semantic names (e.g., ``mass'', ``acceleration''), the AI Scientist completely ignores empirical evidence from a counterfactual law ($F = m \cdot a^2$) and hallucinates the standard textbook law ($F = ma$), fabricating statistical validation ($R^2 = 1$). However, when variables are anonymized, the system avoids confident hallucinations but struggles to discover the true law. Our findings expose a critical failure mode in automated science, demonstrating that AI scientists cannot be fully trusted with empirical discovery unless strictly isolated from their pre-trained semantic priors.